\documentclass{beamer}

\usetheme{Madrid}
\usepackage{amsmath,amssymb}
\usepackage{hyperref}
\usepackage{listings}
\usepackage{xcolor}

% ---------------------------
% Listings (Python)
% ---------------------------
\lstset{
  language=Python,
  basicstyle=\ttfamily\small,
  keywordstyle=\color{blue},
  commentstyle=\color{gray},
  stringstyle=\color{purple},
  showstringspaces=false,
  breaklines=true,
  frame=single,
  columns=fullflexible
}

\title{Running Qiskit on IBM Quantum Hardware}
\author{Instructions for FYS5419/9419}
\date{Spring 2026}

\begin{document}

%-------------------------------------------------
\begin{frame}
\titlepage
\end{frame}

%-------------------------------------------------
\begin{frame}{What you will do}
\begin{enumerate}
  \item Create an IBM Quantum Platform / IBM Cloud account and select an access plan
  \item Initialize Qiskit Runtime credentials locally
  \item Discover available backends (QPUs and simulators)
  \item Run circuits using Qiskit Runtime primitives (Sampler / Estimator)
  \item Monitor jobs and retrieve results
\end{enumerate}

\vspace{0.4em}
\scriptsize
References: IBM Quantum Platform Quickstart and guides on account setup, credentials, and execution primitives. \\
(\href{https://quantum.cloud.ibm.com/docs/guides/quick-start}{Quickstart},
\href{https://quantum.cloud.ibm.com/docs/guides/save-credentials}{Save credentials},
\href{https://quantum.cloud.ibm.com/docs/guides/execute-on-hardware}{Execute on hardware},
\href{https://quantum.cloud.ibm.com/docs/guides/get-started-with-primitives}{Get started with primitives}).%
\end{frame}

%=================================================
\section{Accounts and access}

\begin{frame}{Step 1: Create an account and choose a plan}
\begin{itemize}
  \item Create an account on the \textbf{IBM Quantum Platform} / \textbf{IBM Cloud}.
  \item IBM Quantum Platform offers an \textbf{Open (free)} access path and paid plans (depending on region/org).%
  \item For QPU jobs through \textbf{Qiskit Runtime}, you typically need access to a Runtime instance (linked to an IBM Cloud account).
\end{itemize}

\vspace{0.4em}
\scriptsize
IBM docs: account setup and plan overview.\\
\href{https://quantum.cloud.ibm.com/docs/guides/cloud-setup}{Set up IBM Cloud account},
\href{https://quantum.cloud.ibm.com/docs/guides/plans-overview}{Plans overview},
\href{https://quantum.cloud.ibm.com/registration}{Platform registration}.%
\end{frame}

\begin{frame}{Step 2: Create / select a Qiskit Runtime service instance}
\begin{itemize}
  \item Log in to IBM Quantum Platform and ensure the correct \textbf{account} and \textbf{region} are selected.
  \item Verify you have a \textbf{Qiskit Runtime service instance}; if not, create one.
  \item You will later connect to this instance from Python.
\end{itemize}

\vspace{0.4em}
\scriptsize
IBM docs: initializing the Runtime service account / instance.\\
\href{https://quantum.cloud.ibm.com/docs/guides/initialize-account}{Initialize your Qiskit Runtime service account}.%
\end{frame}

%=================================================
\section{Local environment setup}

\begin{frame}[fragile]{Step 3: Install Qiskit and IBM Runtime packages}
\begin{itemize}
  \item Use a virtual environment (recommended) and install Qiskit + IBM Runtime tooling.
\end{itemize}

\begin{lstlisting}[language=bash]
python3 -m venv .venv
source .venv/bin/activate

pip install --upgrade pip
pip install qiskit qiskit-ibm-runtime qiskit-ibm-provider
\end{lstlisting}

\vspace{0.4em}
\scriptsize
IBM quickstart emphasizes using a virtual environment.\\
\href{https://quantum.cloud.ibm.com/docs/guides/quick-start}{Quickstart}.%
\end{frame}

\begin{frame}[fragile]{Step 4: Save credentials locally (token / API key)}
\begin{itemize}
  \item In IBM Quantum Platform, obtain your API credentials (token/API key).
  \item Save them in your local environment so you do not hard-code secrets in scripts.
\end{itemize}

\begin{lstlisting}
from qiskit_ibm_runtime import QiskitRuntimeService

# Option A (common): save for later use (interactive)
QiskitRuntimeService.save_account(
    channel="ibm_quantum_platform",
    token="YOUR_TOKEN_HERE",
    # instance may be required depending on your plan/org:
    # instance="ibm-q/open/main"   # example format; use your actual instance
    overwrite=True
)

# Later in scripts/notebooks:
service = QiskitRuntimeService(channel="ibm_quantum_platform")
\end{lstlisting}

\vspace{0.4em}
\scriptsize
IBM docs: saving credentials and initializing the service.\\
\href{https://quantum.cloud.ibm.com/docs/guides/save-credentials}{Save your login credentials},
\href{https://quantum.cloud.ibm.com/docs/guides/quick-start}{Quickstart}.%
\end{frame}

%=================================================
\section{Discovering hardware backends}

\begin{frame}[fragile]{Step 5: List available backends (QPUs and simulators)}
\begin{lstlisting}
from qiskit_ibm_runtime import QiskitRuntimeService

service = QiskitRuntimeService(channel="ibm_quantum_platform")

# List backends you can access
backends = service.backends()
print("Number of backends:", len(backends))
for b in backends[:10]:
    print(b.name)
\end{lstlisting}

\begin{itemize}
  \item Choose a backend based on availability, qubit count, and queue length.
  \item IBM provides documentation on viewing backend information via the service API.
\end{itemize}

\vspace{0.4em}
\scriptsize
IBM docs: listing backends and QPU information.\\
\href{https://quantum.cloud.ibm.com/docs/guides/get-qpu-information}{Get backend information with Qiskit}.%
\end{frame}

%=================================================
\section{Running jobs with primitives}

\begin{frame}{Why primitives (Sampler / Estimator)?}
\begin{itemize}
  \item IBM recommends running circuits via \textbf{Qiskit Runtime primitives}:
    \begin{itemize}
      \item \textbf{Sampler}: returns samples / counts from circuit measurements.
      \item \textbf{Estimator}: returns expectation values of observables.
    \end{itemize}
  \item Primitives can include runtime features such as optimized execution and (optionally) error mitigation.
\end{itemize}

\vspace{0.4em}
\scriptsize
IBM docs: primitives and executing on hardware.\\
\href{https://quantum.cloud.ibm.com/docs/guides/execute-on-hardware}{Execute on target hardware},
\href{https://quantum.cloud.ibm.com/docs/guides/primitives-examples}{Primitives examples},
\href{https://qiskit.qotlabs.org/docs/guides/primitives}{Introduction to primitives}.%
\end{frame}

\begin{frame}[fragile]{Example A (Sampler): run a Bell circuit on hardware}
\begin{lstlisting}
from qiskit import QuantumCircuit, transpile
from qiskit_ibm_runtime import QiskitRuntimeService, SamplerV2 as Sampler

service = QiskitRuntimeService(channel="ibm_quantum_platform")

# Pick a backend (replace with one you have access to)
backend = service.backend("ibm_brisbane")  # example name

# Build circuit
qc = QuantumCircuit(2, 2)
qc.h(0)
qc.cx(0, 1)
qc.measure([0, 1], [0, 1])

# Transpile to the backend's ISA
isa_circuit = transpile(qc, backend=backend, optimization_level=1)

# Run with the Sampler primitive
sampler = Sampler(backend=backend)
job = sampler.run([isa_circuit], shots=4000)

print("Job ID:", job.job_id())
result = job.result()
print(result)
\end{lstlisting}

\vspace{0.2em}
\scriptsize
IBM docs: execution on hardware uses Sampler/Estimator primitives.\\
\href{https://quantum.cloud.ibm.com/docs/guides/execute-on-hardware}{Execute on hardware},
\href{https://quantum.cloud.ibm.com/docs/guides/get-started-with-primitives}{Get started with primitives}.%
\end{frame}

\begin{frame}[fragile]{Example B (Estimator): expectation value of an observable}
\begin{itemize}
  \item Estimator returns $\langle O\rangle$ for an observable $O$ on the circuit output state.
  \item Typical use: VQE, QML, metrology objectives, etc.
\end{itemize}

\begin{lstlisting}
from qiskit import QuantumCircuit, transpile
from qiskit.quantum_info import SparsePauliOp
from qiskit_ibm_runtime import QiskitRuntimeService, EstimatorV2 as Estimator

service = QiskitRuntimeService(channel="ibm_quantum_platform")
backend = service.backend("ibm_brisbane")  # example

# Prepare |Phi+> and estimate <ZZ>
qc = QuantumCircuit(2)
qc.h(0); qc.cx(0, 1)

obs = SparsePauliOp.from_list([("ZZ", 1.0)])

isa_circuit = transpile(qc, backend=backend, optimization_level=1)

estimator = Estimator(backend=backend)
job = estimator.run([(isa_circuit, obs)])
print("Job ID:", job.job_id())
res = job.result()
print(res)
\end{lstlisting}

\vspace{0.2em}
\scriptsize
IBM docs: Estimator and examples.\\
\href{https://quantum.cloud.ibm.com/docs/guides/primitives-examples}{Primitives examples},
\href{https://quantum.cloud.ibm.com/docs/guides/get-started-with-primitives}{Get started with primitives}.%
\end{frame}

%=================================================
\section{Sessions, monitoring, and practical advice}

\begin{frame}[fragile]{Sessions (recommended for multiple runs)}
\begin{itemize}
  \item If you execute \emph{many} primitive calls, using a session can reduce overhead.
  \item A session keeps context on the backend for a window of time (useful for iterative workflows).
\end{itemize}

\begin{lstlisting}
from qiskit_ibm_runtime import QiskitRuntimeService, Session
from qiskit_ibm_runtime import SamplerV2 as Sampler

service = QiskitRuntimeService(channel="ibm_quantum_platform")
backend = service.backend("ibm_brisbane")  # example

with Session(service=service, backend=backend) as session:
    sampler = Sampler(session=session)
    job = sampler.run([isa_circuit], shots=2000)
    print(job.result())
\end{lstlisting}

\vspace{0.2em}
\scriptsize
IBM docs: primitives + execution flow; sessions are part of Runtime usage patterns.\\
\href{https://quantum.cloud.ibm.com/docs/guides/execute-on-hardware}{Execute on hardware}.%
\end{frame}

\begin{frame}[fragile]{Monitoring jobs and handling queues}
\begin{lstlisting}
# For any job returned by Sampler/Estimator:
print("Status:", job.status())

# Blocking wait (simple):
result = job.result()

# Some jobs expose logs/metadata depending on runtime configuration:
# print(job.metrics())  # if available in your environment
\end{lstlisting}

\begin{itemize}
  \item Hardware runs are queued; queue times vary by backend load and access plan.
  \item Start with simulators, then validate on hardware.
  \item Keep shot counts modest during development; increase shots for final statistics.
\end{itemize}

\vspace{0.2em}
\scriptsize
IBM docs: backend info and execution guidance.\\
\href{https://quantum.cloud.ibm.com/docs/guides/get-qpu-information}{Get QPU information},
\href{https://quantum.cloud.ibm.com/docs/guides/execute-on-hardware}{Execute on hardware}.%
\end{frame}

\begin{frame}{Practical tips for running on IBM QPUs}
\begin{itemize}
  \item \textbf{Transpile matters:} always transpile to the backend target/ISA.
  \item \textbf{Noise:} expect deviations from ideal counts/expectations; compare with a simulator baseline.
  \item \textbf{Measurement error mitigation:} primitives may support resilience options depending on plan.
  \item \textbf{Reproducibility:} record backend name, calibration date/time, transpiler settings, and shots.
  \item \textbf{Iteration:} develop locally (simulator), then run a small calibration set on hardware.
\end{itemize}

\vspace{0.4em}
\scriptsize
IBM docs: primitives can provide sophisticated implementations including mitigation.\\
\href{https://qiskit.qotlabs.org/docs/guides/primitives}{Introduction to primitives}.%
\end{frame}

%=================================================
\section{Minimal ``Hello World'' checklist}

\begin{frame}{Checklist (quick workflow)}
\begin{enumerate}
  \item Create IBM Quantum Platform / IBM Cloud account and select an access plan.
  \item Create/select a Qiskit Runtime instance for your account/region.
  \item Install packages: \texttt{qiskit}, \texttt{qiskit-ibm-runtime}, \texttt{qiskit-ibm-provider}.
  \item Save credentials with \texttt{QiskitRuntimeService.save\_account(...)}.
  \item List backends: \texttt{service.backends()}.
  \item Build circuit, transpile to backend, run via \texttt{Sampler} or \texttt{Estimator}.
  \item Retrieve results and document settings.
\end{enumerate}

\vspace{0.4em}
\scriptsize
IBM docs hub: \href{https://quantum.cloud.ibm.com/docs/guides/quick-start}{Quickstart}.%
\end{frame}

\end{document}
